\begin{proposition}[counterexample, with corrected family]\hypertarget{tag-01QD}{}\label{fam:src:B.2.3}\label{fam:st:B.2.3}
\StatementRecord{fam:src:B.2.3}{proved}{01QD}{\SeriesRef{fam:src:A.3.1}, \SeriesRef{fam:src:A.2.1}, \SeriesRef{fam:src:B.2.1}, \SeriesRef{fam:src:A.4.1}; \cite[Remark 8.1]{Engel2026}}

The assertion in \cite[Remark 8.1]{Engel2026} of a global proper family over the original punctured $u$-disc for fixed sphere logarithmic data requires a double-cover qualification. The retained family with such data does not descend through $u=w^2$. It does give a proper holomorphic family on each simply connected $u$-open set, and a global proper family over the punctured $w$-disc as in \ref{fam:src:A.3.1}.
\end{proposition}


\begin{proof}
\begin{enumerate}[label=\arabic*.,leftmargin=*]
\item Descent through $u=w^2$ would furnish an isomorphism between the two pulled-back fibres at $w$ and $-w$, namely $X_c$ and $X_{c+1}$. \ref{fam:src:A.2.1} excludes that isomorphism by the nonzero class of $d$ in $\Lambda^{A_2}/N_2\Lambda$. This is an exact order-two obstruction on the normalized finite cover. It applies to all sphere data by \ref{fam:src:B.2.1} because $r_2$ is odd.

\item There is also a direct analytic witness near zero. The relative Albanese invariant of the multiplicity-four reduced fibre in a holomorphic descended family would be a single-valued holomorphic function of $u$. After pullback it would be even in $w$. But \ref{fam:src:A.4.1} gives $j_2(w)=C^{-1}w^{-1}+744+O(w)$, where $C=e^{-\pi i\kappa_2^{\rm part}}\ne0$; its Laurent coefficient of $w^{-1}$ is nonzero. Thus it is not even on any punctured neighbourhood. This verifies the obstruction independently of a choice of isomorphism lift. Local inverse logarithms of $u$ give the proper local families already constructed; coherent even shifts give their global double-cover family. The untwisted middle quotient itself has the integer shear identification, so this refutation concerns the fixed sphere logarithmic data in the stated remark.
\end{enumerate}
\end{proof}